\chapter{Anforderungen und Konzeption des Scheduling-Visualizers}

Dieses Kapitel definiert die fachlichen, didaktischen und technischen Anforderungen der Anwendung und leitet daraus das konzeptionelle Loesungsdesign ab. Im Mittelpunkt steht die Frage, wie präemptive Scheduling-Verfahren fachlich korrekt, didaktisch nachvollziehbar und mit hoher Interaktionsqualitaet visualisiert werden koennen. Die konzeptionellen Entscheidungen orientieren sich dabei an etablierten Scheduling-Grundlagen \cite{silberschatz2018operating,tanenbaum2014modern,pinedo2016scheduling}, an Erkenntnissen aus der Algorithmusvisualisierung \cite{hundhausen2002meta,naps2002engagement} sowie an Anforderungen der Usability und Accessibility \cite{brooke1996sus,wcag2018}.

\section{Zielgruppe und Einsatzkontext}

Primaere Zielgruppe sind Studierende in einführenden und vertiefenden Lehrveranstaltungen zu Betriebssystemen. Die Anwendung soll insbesondere in Situationen unterstuetzen, in denen formale Algorithmusbeschreibungen allein nicht ausreichen, um den zeitlichen Ablauf präemptiver Entscheidungen zu verstehen.

Der Einsatzkontext umfasst sowohl Lehrveranstaltungen mit gefuehrter Demonstration als auch das selbststaendige Lernen. Daraus folgen zwei zentrale Rahmenbedingungen: Erstens muss die Anwendung reproduzierbare Ergebnisse liefern, damit Beobachtungen zwischen verschiedenen Nutzenden vergleichbar bleiben. Zweitens braucht sie eine flexible Interaktion, damit sowohl exploratives Ausprobieren als auch gezieltes Nachstellen vorgegebener Szenarien moeglich ist.

\section{Fachliche Anforderungen}

Aus fachlicher Sicht muss der Visualizer die betrachteten Verfahren konsistent und korrekt modellieren. Dazu gehoeren insbesondere LCFS, Round Robin, Strict Priority und MLFQ. Relevante Kernaspekte sind Ankunftszeit, Burst-Zeit, Prioritaet, Zeitquantum, Präemption und Kontextwechsel \cite{silberschatz2018operating,tanenbaum2014modern}.

Zusatzlich muessen die wichtigsten Vergleichsmetriken nachvollziehbar berechnet werden. Dazu zaehlen Wartezeit, Durchlaufzeit, Reaktionszeit, Kontextwechselhaeufigkeit und weitere Kenngroessen zur Bewertung von Fairness und Ressourcennutzung \cite{pinedo2016scheduling}. Die konzeptionelle Anforderung lautet daher, dass Simulationslogik und Metrikberechnung auf derselben Zustandsbasis arbeiten und keine voneinander getrennten Berechnungspfade verwenden.

\section{Didaktische Anforderungen}

Didaktisch soll die Anwendung nicht nur Ergebnisse anzeigen, sondern den Entscheidungsprozess der Algorithmen sichtbar machen. Lernende sollen erkennen koennen, warum ein Prozess zu einem bestimmten Zeitpunkt ausgewaehlt, unterbrochen oder erneut eingeplant wird. Dieser Anspruch folgt direkt aus der Forschung zu lernwirksamer Algorithmusvisualisierung, die insbesondere aktive Interaktion und vergleichende Betrachtung als wirksame Lernhebel beschreibt \cite{hundhausen2002meta,naps2002engagement,stasko1998software}.

Die Konzeption priorisiert deshalb drei didaktische Leitprinzipien:

\begin{itemize}
    \item \textbf{Transparenz der Zustandsuebergaenge:} Ereignisse wie Start, Präemption und Abschluss muessen nachvollziehbar dargestellt werden.
    \item \textbf{Vergleichbarkeit:} Mehrere Algorithmen sollen unter identischen Eingabedaten gegenuebergestellt werden koennen.
    \item \textbf{Kontrollierbares Lerntempo:} Schrittweises und automatisches Abspielen sollen gleichwertig unterstuetzt werden.
\end{itemize}

\section{Funktionale Anforderungen}

Funktional muss die Anwendung den gesamten Ablauf von der Szenariodefinition bis zur Ergebnisinterpretation abdecken. Daraus ergeben sich folgende Kernfunktionen:

\begin{itemize}
    \item Erfassung und Bearbeitung von Prozessmengen mit relevanten Attributen.
    \item Auswahl eines Scheduling-Verfahrens und seiner Parameter.
    \item Start, Pause, Schrittmodus und Zeitsprung innerhalb der Simulation.
    \item Visualisierung der Ausfuehrungssegmente auf einer Zeitachse.
    \item Ausgabe vergleichbarer Metriken pro Lauf.
    \item Gegenueberstellung mehrerer Laeufe auf identischem Szenario.
\end{itemize}

Konzeptionell wird damit ein durchgaengiger Workflow vom Eingabemodell zur Auswertung abgebildet. Jede Funktion muss auf denselben Simulationsdaten aufbauen, um Inkonsistenzen zwischen Anzeige und Metrikberechnung zu vermeiden.

\section{Nicht-funktionale Anforderungen}

Nicht-funktionale Anforderungen betreffen vor allem Nachvollziehbarkeit, Robustheit und Nutzbarkeit. Die Anwendung soll deterministisch arbeiten, damit identische Eingaben zu identischen Ergebnissen fuehren. Dies ist fuer didaktische Vergleichbarkeit und wissenschaftliche Reproduzierbarkeit zentral \cite{banks2010discrete}.

Darueber hinaus gelten Anforderungen an Performance und Bedienqualitaet: Interaktionen muessen unmittelbar rueckmelden, die Darstellung muss auch bei groesseren Szenarien stabil bleiben, und zentrale Funktionen muessen ohne Maus nutzbar sein. Anforderungen an semantische Benennung, Fokusfuehrung und Kontrast orientieren sich an etablierten Accessibility-Kriterien \cite{wcag2018}. Fuer die spaetere Nutzerbewertung wird ein standardisiertes Usability-Instrument wie SUS als geeignet betrachtet \cite{brooke1996sus}.

\section{Konzeption der Simulationslogik}

Die Simulationslogik wird als deterministischer Kern konzipiert, der diskrete Zeitschritte verarbeitet und Ereignisse in wohldefinierter Reihenfolge erzeugt. Die zentrale Entscheidung besteht darin, Simulationszustand, Ereignisprotokoll, Zeitachsen-Segmente und Metriken gemeinsam aus einem einheitlichen Zustandsmodell abzuleiten. Dadurch bleiben alle Ausgaben konsistent.

Konzeptionell werden dabei zwei Anforderungen verbunden: erstens fachliche Korrektheit der Scheduling-Regeln, zweitens auswertbare Verlaufsdaten fuer Visualisierung und Evaluation. Das Ereignismodell muss daher nicht nur Endergebnisse liefern, sondern auch Zwischenschritte explizit protokollieren, damit Präemptionen und Kontextwechsel analytisch nachvollzogen werden koennen \cite{silberschatz2018operating,pinedo2016scheduling}.

\section{Konzeption der Visualisierung}

Als primaeres Visualisierungsmedium wird ein Gantt-artiger Zeitverlauf eingesetzt. Diese Form ist geeignet, da sie parallel Informationen ueber Prozessdauer, Unterbrechungen und Wartephasen abbildet und damit den Vergleich unterschiedlicher Strategien erleichtert.

Die Visualisierung wird konzeptionell als Ableitung des Simulationskerns verstanden, nicht als eigene Logikinstanz. Das bedeutet: Die Darstellung interpretiert vorhandene Zustands- und Segmentdaten, trifft aber keine eigenstaendigen Scheduling-Entscheidungen. Diese Trennung reduziert Fehlerquellen und erhoeht die Validierbarkeit. Fuer die Interaktion werden ergänzend Tooltips, Hervorhebungen und Zeitsteuerung vorgesehen, um sowohl globale als auch detailorientierte Analyse zu ermoeglichen \cite{hundhausen2002meta,stasko1998software}.

\section{Konzeption der Benutzungsoberflaeche}

Die Benutzungsoberflaeche wird in drei logisch getrennte Bereiche gegliedert: Eingabe und Konfiguration, Simulationssteuerung sowie Ergebnisdarstellung. Diese Struktur soll die mentale Last reduzieren und den Wechsel zwischen Experimentieren und Interpretieren erleichtern.

Zentrale Interaktionsprinzipien sind Konsistenz, Sichtbarkeit des Systemzustands und fehlertolerante Bedienung. Aktionen muessen eindeutig rueckgemeldet werden, insbesondere beim Wechsel zwischen pausiertem und laufendem Zustand. Modale Dialoge und Menues sind mit klarem Fokusverhalten und semantischer Beschriftung auszustatten, damit die Anwendung auch mit assistiven Technologien nutzbar bleibt \cite{wcag2018}. Insgesamt bildet die Oberflaechenkonzeption damit die operative Umsetzung der in Kapitel 3 und Kapitel 4 abgeleiteten didaktischen und HCI-Anforderungen.
